\newpage\thispagestyle{empty}
\setcounter{footnote}{0}

\chapter*{Acknowledgements}

%\begin{quote}
My sincere thanks to Enxin Wu for inviting me to give a series of lectures on diffeology at Shantou University,
as part of a ``Guangdong Science and Technology Project''.%
\footnote{\ \begin{CJK}{UTF8}{gbsn}广东省科技计划项目 合同书 (“海外名师”项目) 项目编号: 2020A1414010269.\end{CJK}}
This led to the publication of these lecture notes,
giving me the opportunity to present in China this new look at differential geometry,
and to show some examples of applications.

\medskip\noindent
It is always a pleasure for me to thank the Hebrew University of Jerusalem, Israel,
to which I belong by heart.
I am thinking of Professor Elon Lindenstrauss and Professor Jay Fineberg in particular,
who a few years ago,
as Chairman of the Mathematics Department and Dean of the Faculty of Science at the time,
enabled me to continue my academic activity as a visiting professor,
despite my retirement from the CNRS in France.
For that,
I am deeply grateful.
It is thanks to them and to all my colleagues at HUJI,
Emanuel Farjoun,
Itamar Cwik,
Yael Karshon,
Jake Solomon,
to cite only a few,
and all those who have always welcomed me,
that this book exists.

\medskip\noindent
I also would like to take this opportunity to thank my Japanese colleagues,
in particular Norio Iwase and Katsuhiko Kuribayashi,
for their many invitations to participate in the conferences on differential homotopy,
where diffeology is also abundantly discussed,
that they have been organizing regularly in Japan for several years.
They are a source of inspiration and questioning,
as well as prospects for further development.
I hope that these notes will be useful to their students,
some of which were inspired by discussions I have had with them.

\medskip\noindent
And lastly,
I would like to express my sincere gratitude and thanks,
for the second time,
to all the members of the Beijing WPC team who have helped me so much,
in particular to the publishing director Liang Chen,
and especially to Yeqing Liu,
my ever-patient editor,
who made this publication possible,
put the finishing touches on it,
and with whom I am preparing a new book,
this time on mechanics: "The Geometry of \{Motion\}".

\medskip\noindent
Finally,
I wish to thank DeepSeek (\begin{CJK}{UTF8}{gbsn}深度求索\end{CJK}, \url{www.deepseek.com}),
the AI assistant that accompanied me throughout the revision of these lectures.
I am grateful to its developers for creating such a powerful tool,
and to the company for their decision to keep it freely accessible,
allowing researchers to work uninterrupted over long stretches of time,
dedicated to the advancement of science.
Mathematics is not a series of isolated queries but a sustained,
immersive exploration.
DeepSeek helped me correct typos,
resolve notational ambiguities,
and refine the English prose.
And more than a mere proofreader,
it was also a true collaborator,
catching inconsistencies,
suggesting improvements,
and engaging in the mathematical substance of the work.
Its assistance made this revised edition a better book.
%\end{quote}

\vfill
\begin{flushright}
  Patrick Iglesias-Zemmour\\
  Einstein Institute of Mathematics \\
  Edmond J. Safra Campus \\
  The Hebrew University of Jerusalem \\
  Givat Ram. Jerusalem, 9190401, Israel \\
  \medskip
  \url{http://math.huji.ac.il/~piz}\\
  \url{piz@math.huji.ac.il}\\
  \medskip
  \hfill December 18, 2023 \\
  \hfill Revised January 2026
\end{flushright}
